\section*{Computer Algebra Systems: source notes}
\begin{enumerate}
\item See the exact source-note record in the archive.
\end{enumerate}
\section*{Input-Output Analysis: source notes}
\begin{enumerate}
\item All three supplied answers to Exercise 1 shift the already-adjusted next-year example, rather than the current-year external demands stated in the questions. Each steel RHS is 200 too high and each auto RHS is 25 too low for its prompt.
\end{enumerate}
\section*{Accuracy of Computations: source notes}
\begin{enumerate}
\item A semicolon immediately terminates the pivot-comparison if statement. The following max assignment is unconditional and chooses the last visited row, not necessarily the largest absolute entry.
\item For the stated 253 + 2/3, the answer writes .25 times 10\textasciicircum{}2 = 25. Under its intended two-significant-digit restriction the rounded sum is 250, written .25 times 10\textasciicircum{}3. The prompt says decimal places while the answer uses significant digits; those are not identical conventions.
\item Exercise 3 starts with right-hand side 1.559, but the supplied rounded reduction changes the unchanged first row to 1.569. This is not rounding to four significant digits. The source is preserved; the complete four-digit worked answer has not been repaired here.
\item The answer derives the unique solution (1,1,1) by division. Uniqueness requires epsilon not equal to 0 or 5/9. At either exceptional value (1,1,1) is still a solution, but the displayed system has infinitely many solutions.
\item Exercise 5(b) starts from a matrix that is not the preceding system with epsilon = .001: its row-2 and row-3 x coefficients change, and four epsilon-dependent coefficients are ten times too small. Later rounded steps must not be treated as a verified calculation for the stated system.
\end{enumerate}
\section*{Analyzing Networks: source notes}
\begin{enumerate}
\item The traffic-circle question gives cars per ten minutes, while answer (e) says five minutes. Both passages are preserved; use the question’s ten-minute unit when interpreting the data.
\item Exercise 5’s prose states i1+25=i2+125, implying i1−i2=100. Its immediately following system instead uses i1−i2=10. The table and labelled intersection give i1+65=i2+75, supporting the system’s 10, not the prose’s 100. Original text and system remain unchanged.
\item The Wheatstone-bridge derivation divides by branch currents and resistances. The stated resistance ratio uses the physical assumptions of nonzero excitation and positive resistances; if every current is zero, zero detector current alone does not determine that ratio.
\end{enumerate}
\section*{Vector Spaces: source notes}
\begin{enumerate}
\item In the proof of scalar identity for matrices, the final matrix contains sa, sb, sc, sd. Those entries must be a, b, c, d: the scalar is 1, not a free s.
\item In the closure-under-addition display for upper-triangular matrices, an equals sign is missing between the second addend and the resulting sum matrix.
\item For the operation r·(x,y)=(rx,0), the preceding distributivity counterexample does not work: both sides equal (0,0). The separately stated scalar-identity counterexample is valid and the answer “not a vector space” remains correct.
\item The requested arrowed-operation restatement is inconsistent in axioms (6)–(8): scalar multiplication still uses plain · in (6)–(7), and vector addition uses plain + in (8). These should use the newly defined arrowed scalar and vector operations; scalar addition r+s remains plain.
\item The span-membership row reduction omits a step. The two labelled row operations produce third row (0, 3/2, 7/2), not (0,0,3). Adding the resulting second row to the third then gives (0,0,3), so the “no solution” conclusion is unchanged.
\item In the three-dimensional superhero answer, the three device vectors span all of R³, not a plane. Their column matrix has determinant 18. The “no place to hide” conclusion and the displayed elimination are correct.
\item The span-extension proof should assume v belongs to span(S), not necessarily S. Its displayed linear-combination representation is exactly the weaker hypothesis needed by the proof.
\item The intersection proof changes its vector names: it assumes w and s belong to the intersection, then forms r v+s w. Use v and w consistently as vectors, with r and s as scalars.
\item The union-of-subspaces counterexample names ambient space R³ but displays two-entry coordinate vectors. Either use R² as the ambient space or append a zero third coordinate to each vector.
\item The condition given for span(S∪T)=span(S)∪span(T) is too strong. Equality holds exactly when one span contains the other, not only when one original set contains the other. For S=\{(1,0)\} and T=\{(2,0)\}, neither set contains the other but both spans are the same line. The following noncontainment argument confuses membership in a set with membership in its span.
\end{enumerate}
\section*{Linear Independence: source notes}
\begin{enumerate}
\item In the \{1+x,1-x\} example, R2 - R1 yields -2c2 = 0, not +2c2 = 0. The conclusion of independence is unchanged.
\item The general echelon-matrix statement needs the qualification nonzero rows. An echelon matrix with a zero row supplies an immediate counterexample. The displayed example itself has no zero row.
\item When isolating v from 0 = c1 s1 + ... + cn sn + c(n+1) v, each displayed coefficient of si needs a minus sign. The span-membership conclusion remains valid.
\item The first equation in the spanning-set reduction contains a plus sign in the empty c4 column, printing two successive plus signs. The intended first equation is c1 + c3 + 3c5 = 0.
\item In the polynomial exercise answer, the stated row operations give the last pivot -152/13, not -128/13. The independence conclusion remains correct.
\item In the three-row-vector exercise, the operation producing the zero third row must be R3 - (4/7)R2, not R3 + (4/7)R2. The final dependence vector is correct.
\item The union example defines S and T but its last sentence calls them S1 and S2. These are the same displayed sets; their union is dependent.
\item The question labels its proposed direction as if, but the proposed implication from an independent union to trivial intersection is the only-if direction. This note corrects the direction label only; it supplies no answer.
\item The supplied answer reverses the names if and only if: trivial intersection implies an independent union is the if direction; the reverse is only if. The corrected criterion and proof substance are unchanged.
\item The two-by-two determinant argument needs unique solution, not solution; every homogeneous system has a zero solution. Its combined fraction also needs denominator a, not d.
\item The three-by-three elimination step needs the negative multiplier -((ah-bg)/a) on row 2 to clear row 3. The displayed final pivot is the result of this subtraction.
\item In the a != 0, ae-bd = 0 case, the proof describes singularity while calling it nonsingularity. After the row swap the system is nonsingular exactly when both ah-bg and af-cd are nonzero, equivalently when the determinant is nonzero. Either factor zero, and the subsequent equal-zero calculation, establish singularity.
\item For the two vectors in R3, dependence requires ae-bd = ah-gb = dh-ge = 0. Equality of these three minors without the final zero condition is insufficient.
\end{enumerate}
\section*{Basis and Dimension: source notes}
\begin{enumerate}
\item The first polynomial-basis answer omits the contributions of c1 to the x and constant coefficients. Its proposed c2 and c3 therefore do not represent a general quadratic when a2 is nonzero. The stated set is still a basis. Correct coefficients (c1,c2,c3) are (a2; a0/2 + a1/4 - a2/4; -a0/2 + a1/4 + 3*a2/4).
\item The matrix-basis proof obtains c3=0 from the lower-right entries, not the lower-left entries as stated at line 918. The displayed third basis matrix has its sole nonzero entry at the lower right. The conclusion remains correct.
\item Choose s2 in S outside span(B1), and repeat that choice within S. Such an element exists whenever span(B1) is a proper subspace of span(S). B1 is initially a basis of its own span, not necessarily the full space.
\item Read the intermediate set as \{-1+x, 2x\textasciicircum{}2+x\textasciicircum{}3\}. The printed x\textasciicircum{}2+x\textasciicircum{}3 fails the exercise condition a2-2a3=0; the final displayed basis and dimension are correct.
\item The two displayed sets describe the same diagonal-matrix space; an equality sign is missing at the start of the second line. Preserve the original display and identify this separately.
\item Use a nontrivial relation among the concatenated bases. Independence within each basis implies that the partial combination from either basis is nonzero. These equal, opposite-signed partial combinations provide a nonzero vector in U intersect W; no individual basis vector need lie in the other space.
\item Read V=W in the first sentence of this part as U=W. The remainder of the argument already uses U and W.
\item The right-hand column in the first displayed equation should be (1,3), as in the question and subsequent system. The conclusion that it is outside the column space remains correct.
\item At the second arrow, read the row-four operation as -2 times row one plus row four. The displayed echelon matrix then follows. The later basis ending in -5x\textasciicircum{}2 is valid: scaling the nonzero third basis vector does not change its span.
\item Read “at least two” here as “at most two”. The fourth and fifth columns are independent, so the column rank is exactly two, as stated.
\item For unequal row counts, first append zero rows to the shorter matrix, or compare only the nonzero reduced rows. The exercise claims equality of nonzero rows, not equality of differently sized matrices. A proof can stack A over B and B over A: the two stacked matrices are row equivalent, and reducing either dependent extra block leaves the same nonzero rows.
\item In part (b), read the right-hand entries as d1 through dm. There is one right-hand entry for each of the m equations.
\item For the missing direction, if every right-hand vector in R\textasciicircum{}m is attainable, the columns span R\textasciicircum{}m. The column rank is then m, and equality of row and column ranks gives full row rank. The second printed paragraph is the contrapositive of the first direction.
\item In part (e), W1 is a subspace: shifting its parameter still describes the entire x-axis. W2 is not a subspace because it does not contain zero. The answer “No” remains correct.
\item Read “greater than” here as “at least”. Dimensions eight, nine and ten for the sum, and eight, seven and six for the intersection, are all possible. The printed final list is correct.
\item Read “every subset” as “every subspace”. Closure under addition is available because W is a subspace; the identity W+W=W is correct.
\end{enumerate}
\section*{Fields: source notes}
\begin{enumerate}
\item The displayed axiom list does not explicitly require the additive and multiplicative identities to be distinct. Its literal conditions also hold in the singleton algebra with 0=1; the usual field definition excludes this case. Record the missing explicit condition separately, without changing the pinned source.
\item The supplied binary-field answer calls addition closure and commutativity the first half of condition (2), but those appear in condition (1). The operations and the finite-field conclusion are correct.
\end{enumerate}
\section*{Crystals: source notes}
\begin{enumerate}
\item Exercise 1’s supplied answer uses 3.34 × 10⁻¹⁰ cm, but the topic’s own 3.34 Å and Å-to-metre definition give 3.34 × 10⁻⁸ cm. It also divides one edge length even though the question asks for regions in a square face; using the printed spacing, the face count is (0.1/(3.34 × 10⁻⁸))², approximately 8.96 × 10¹². This checks the internal conversion and area count, not the physical lattice constant.
\item Exercise 2(a)’s displayed system gives c₁ = 47447/17466 ≈ 2.72 and c₂ = 314/123 ≈ 2.55, not c₁ ≈ 2.74. Part (b) then reverses the rounded coefficients to 2.55β₁ + 2.72β₂; both original diagram definitions use that reversed pair. The prose and diagrams are retained unchanged rather than silently repaired.
\item Exercise 3 lists (0.25,0.25,0.25) as an atom a quarter of the way down from the top. That point has height 0.25, not 0.75. The original diagram instead explicitly places this atom at (0.25,0.25,0.75), consistent with the prose description. The original supplied answer remains unchanged.
\item Exercise 4(a) prints 195.08/6.02 × 10²³ = 3.239 × 10⁻²². Its denominator needs grouping, and the question gives 6.022 rather than 6.02. Using the stated data gives 195.08/(6.022 × 10²³) ≈ 3.239455 × 10⁻²² grams per atom. The other printed values are retained; no independent physical-constant verification is claimed.
\end{enumerate}
\section*{Voting Paradoxes: source notes}
\begin{enumerate}
\item Exercise 4(c) gives strict triangle inequalities as a necessary-and-sufficient condition, but the topic and Exercise 5 explicitly allow tied pairwise totals. For a=b=1 and c=0 the vote vector is (0,2,0), which meets that stated definition but fails |a|<|b+c|. Under the source’s nonnegative/nonpositive convention, use the three weak inequalities (≤), together with the same-sign condition. The printed strict inequalities characterize a cycle with no tied totals. The original answer is preserved unchanged.
\end{enumerate}
\section*{Dimensional Analysis: source notes}
\begin{enumerate}
\item The pendulum discussion changes g to an undefined a in one product. Its quantity table, preceding product, exponent system and following basis all use g for gravity. The consistent symbol here is g. The original formula remains unchanged.
\item Exercise 1’s answer lists xt/v₀² as dimensionless, but its dimensional formula is L⁻¹T³. The immediately preceding exponent vector (1,0,−2,0,1,0) instead gives xg/v₀², which is dimensionless. Original wording and mathematics are retained.
\item Exercise 4’s supplied reduction labels an operation ρ₁+ρ₄ although the displayed system has only three rows. The resulting row is the sum of the first and third rows, (0,1,1,−1), after dropping the zero row. The original label is preserved; its row index should be 3, not 4.
\item Exercise 5 needs a domain convention for real powers. Arbitrary real bases do not suffice: for the zero homogeneous system with one variable m=−1, exponent 1 gives −1, but multiplying this vector by scalar 1/2 would require a real square root of −1. With m=0, exponent −1 is undefined. One valid interpretation uses formal monomials identified by exponent vectors in the homogeneous kernel; another uses positive bases with well-defined real-power laws. No such restriction is silently added to the original question.
\end{enumerate}
\section*{Maps Between Spaces: source notes}
\begin{enumerate}
\item The displayed injection has codomain R\textasciicircum{}3 but its image is only the xy-plane. It is an isomorphism onto that plane, not onto R\textasciicircum{}3. The subsequent proof also switches its name from iota to f.
\item The concluding domain-vector expression in the dependence proof contains stray f symbols and an unmatched parenthesis. The immediately preceding correct display has the domain sum with no f applied to its terms.
\item The final reflection formula loses the numerator/denominator grouping in its inline quotients. It differs from the two correctly grouped fractions derived just above. At slope k=1 it should exchange x and y; the literal printed expression does not.
\item The hint equating “the second column is not a multiple of the first” with nonzero determinant needs the first column to be nonzero. For (a,c)=(0,0) and (b,d)=(1,0), the second column is not a multiple but ad-bc=0. Nonzero first column follows in the automorphism argument, but is absent from the standalone hint.
\item The last line of the coordinate-map linearity calculation omits the scalars c and d. Its preceding line is c Rep\_B(p)+d Rep\_B(q), not Rep\_B(p)+Rep\_B(q).
\item The five displayed vectors do form a basis, but the worked decomposition is wrong: the final two basis vectors both have polynomial component 1, so its right side has constant term 12 instead of 3. For that unchanged basis, the first coefficient should be -6 instead of 3.
\item The injectivity premise for the internal/external direct-sum map omits f on the right side. The next displayed equality of sums makes clear that both ordered pairs must first be mapped by f.
\item The supplied coefficient-dot-product and coefficient-derivative answers assume the monomial basis B=(1,x,x\textasciicircum{}2,x\textasciicircum{}3), which the question does not specify. For arbitrary B the pullback inner product and transported derivative depend on B.
\item The proposed basis correspondence is written f:B to W and called onto, although its finite image D is a basis, not the whole nontrivial real space W. The intended bijection is between B and D, followed by its linear extension from V to W.
\item The triangle diagram ch3.73 labels the line angle theta, while the accompanying derivation uses phi for this same line angle (theta names the input-vector angle in the previous diagram). The source diagram is retained, with this label mismatch disclosed separately.
\end{enumerate}
\section*{Homomorphisms: source notes}
\begin{enumerate}
\item The first supplied linearity calculation uses c\_2 where x\_2 belongs in the second component of the second column vector. With c\_2=2, x\_2=3 and y\_2=z\_2=0, that contribution is printed as 4 instead of 6.
\item The last line of the matrix-functional linearity calculation lacks the closing parenthesis after the first h argument. It belongs immediately after the first matrix, before the next plus sign.
\item The first image vector in this displayed dependence calculation has lost its subscript 1. The indexed-family argument also requires retaining repeated images, as explained in the separate dependence note below.
\item The final transpose calculation keeps a\_\{j,i\} and b\_\{j,i\} inside the newly reintroduced transpose signs. Those inner entries should be the original a\_\{i,j\} and b\_\{i,j\}; otherwise the displayed expression transposes twice.
\item This statement is valid for an indexed family with repeated images retained, but not for the image as an ordinary set. Under h(x,y)=x, the dependent set \{(1,0),(1,1),(1,2)\} has image \{1\}, which is independent. The proof preserves indexed terms; the wording must be read with that distinction. A later supplied exercise answer explicitly notes that sets remove repeated elements.
\item Preserving independence requires injectivity, not surjectivity onto the declared codomain. The map (x,y) to (x,y,0) preserves independence but is not an isomorphism onto R\textasciicircum{}3. It is an isomorphism onto its range, as the following text correctly states.
\item This exercise declares polynomials of degree at most 3, but specifies h only for ax\textasciicircum{}2+bx+c. The action on the cubic term is missing. A degree-at-most-2 domain or a specified cubic image would repair the definition; neither is silently chosen here. In its supplied preimage argument, b=x and c=y also need to be stated.
\item The last denominator of this null-space condition lost its grouping: a\_n/(n+1) is required, not a\_n/n+1. The preceding displayed integral already contains the correct fraction.
\item The onto-map construction needs a separate zero-dimensional-codomain case: when k=0 there is no w\_k. Sending all remaining basis vectors to zero instead works for every k, including 0. The positive-dimensional case as printed is otherwise valid.
\item The rank-one transformation proof temporarily calls the map h although the exercise and subsequent equations call it t. Those two occurrences of h denote the same t, not a newly defined map.
\item The proposed summands have rank at most one, not necessarily rank one: a summand is zero when h(beta\_i)=0. Omit zero summands, with an empty-sum convention for the zero map. The final displayed linearity check also omits the vector factor h(beta\_i): its middle term must be (r c\_i+s d\_i)h(beta\_i), not the scalar r c\_i+s d\_i.
\item The example proves that the onto-homomorphism relation is not symmetric, not that it is non-reflexive. Every vector space maps onto itself by its identity map.
\end{enumerate}
\section*{Computing Linear Maps: source notes}
\begin{enumerate}
\item The introductory prose says the map is applied to the matrix. In the displayed computation the map is applied to a domain vector; the matrix represents the map.
\item The supplied proof initially uses a row of H and n codomain basis vectors. The definition requires column i: h(beta\_i)=sum over j=1..m of h\_(j,i) delta\_j.
\item The supplied linearity calculation omits h on its final basis vector. The final term must be c\_n h(beta\_n), as in the immediately following expanded line.
\item The proposed chain in the triangularization answer does not necessarily end at W as the question requires. The question also indexes its chain to m while asking for every domain basis index i up to n; unequal dimensions require an explicit index convention.
\item One supplied coordinate vector has an unmatched opening parenthesis before its first 1/2 entry.
\item The displayed kernel generator (-1/2,1) is incorrect for G=(1,2;3,6). The source equations x+2y=0 and 3x+6y=0 give (-2,1), and kernel coordinates belong to the domain basis B, not the codomain basis D.
\item The nullity explanation invokes the codomain representation Rep\_D on W. To transfer the dimension of the kernel coordinates back to the kernel in V it needs the domain representation Rep\_B on V.
\item The source labels a vector representation with the pair B,D. For a domain vector the representation is with respect to B alone; the pair B,D is used for a map representation.
\item The source formula (x,y) maps to (x+3y,y), but the geometric description names the y-axis. This is a horizontal shear: the displacement is parallel to the x-axis, with magnitude three times the distance from the x-axis.
\item The supplied linear-functional proof calls an n-by-1 column the matrix for a map R\textasciicircum{}n to R. That matrix must instead be the 1-by-n transpose row, whose product with a coordinate column is the stated dot product.
\item For the displayed ordered basis D of the two-by-two matrix space, the coordinates are a four-by-one column, not the two-by-two matrix itself. The following coordinate calculation uses the correct four entries.
\item The left side of the image calculation omits the transformation t. The right side computes t(1,-1)=(-4,0), not the coordinates of (1,-1).
\item The last entry in the displayed bottom matrix row is indexed h\_(1,n), whereas that row is row m and requires h\_(m,n).
\item An equality sign is missing between the representation of t and its displayed matrix.
\item The compressed projection display changes the map/matrix labels to h and H, although this example uses the projection pi. These generic names are not defined for this projection display.
\end{enumerate}
\section*{Matrix Operations: source notes}
\begin{enumerate}
\item The original prose repeats the word “of”.
\item The regrouped final coefficient of v\_n in the first output coordinate uses g\_(1,1). It should use g\_(1,n), matching the immediately preceding expansion of g(v).
\item The matrix-product entry formula in the multiplication-cost answer has inconsistent final indices. Its last summand should be g\_(i,r) h\_(r,j), not g\_(1,r) h\_(r,1).
\item The three displayed reciprocal diagonal matrices on the right also carry inverse exponents. As written, the equalities are false; the right-hand inverse exponents should be absent. The following scalar expression also contains a stray underscore.
\item The left-inverse calculation displays a 2-by-2 unknown matrix multiplying the given 2-by-3 matrix but sets the product equal to a 3-by-3 identity. A left inverse must be 3-by-2. The following six-unknown system includes e and f, while the prose incorrectly says four unknowns.
\item The question says exactly one of its two assertions is true, but both are false without additional hypotheses. A left inverse of an injective linear map need not be linear away from the image. The supplied answer incorrectly declares the left-inverse assertion true.
\item This answer incorrectly states that every left inverse of a linear map must be linear. The cited two-sided inverse theorem does not justify that claim on points outside the image.
\item The supplied proof of TS=I if and only if ST=I asserts that these are each equivalent to both matrices being nonsingular, which is false, and only proves two necessary conditions. It does not establish the requested equivalence.
\item The answer names the pair of bases as standard bases of dimensions 2 and 2 while the displayed map has domain R\textasciicircum{}3 and codomain R\textasciicircum{}2.
\item The row-operation label writes (a/ad-bc) times row 2. The displayed reduction requires a/(ad-bc), with the whole denominator grouped.
\item The answer has an incomplete phrase: “the j,i entry of is”. The matrix name after “of” is missing.
\item The rank-product proof bounds an image dimension by the dimension of the domain of g, then concludes a bound by rank(g). That step is not justified: domain dimension need not equal rank(g).
\item The answer counts the powers of a matrix as distinct set elements. They need not be distinct: for T=I, the displayed set is just \{I\}. The dimension argument applies to the indexed family of n-squared-plus-one powers, not necessarily a set of that cardinality.
\item The final shifted derivative drops the coefficient n from its last term. It should end with n times a\_n x\textasciicircum{}n, not a\_n x\textasciicircum{}n.
\item The sentence says “is to rescales the rows”; the verb after “to” should be “rescale”.
\item The answer says the displayed matrix swaps rows one and three, but the matrix and product swap rows one and two, as the question requests.
\item The matrix displayed for g does not represent the stated basis action. The action beta\_1 to beta\_2 and beta\_2 to zero has columns (0,1) and (0,0), whereas the printed matrix has columns (0,0) and (1,0). The sentence also contains “if for we use”.
\end{enumerate}
\section*{Change of Basis: source notes}
\begin{enumerate}
\item The second basis vector in the displayed coordinate system uses y\_1 where the definition immediately above gives y\_2. The two occurrences should use y\_2.
\item The next coordinate system again writes the second basis vector with y\_1 in place of its defined y\_2.
\item The instruction says to concatenate the coordinate columns into a basis; the displayed object is the change of basis matrix.
\item The reflection-inverse derivation includes an extra uninverted factor on the left. That side is a matrix times its own inverse, hence the identity, while the right side is the reflection matrix.
\item The phrase “to translates a representation” should read “to translate a representation”.
\item The sentence “The proof tells us what how the bases change” contains an extra word.
\item The left side of this equality uses B\_3 while the coordinate column and the basis defined alongside it use B\_2.
\item The proof explicitly establishes equal rank implies matrix equivalence, but omits the reverse direction of the stated if-and-only-if result.
\item The question has six matrix subparts, but its answer lists only five size/rank entries. The entry for the fourth matrix is omitted, shifting the final two answer labels.
\item The second identity expansion writes (1,0)=-(1,0)+0(0,1). The first coefficient should be +1.
\item The blanket claim that matrix equivalence classes are not closed under addition omits the zero-rank class exception.
\item The general n-dimensional similarity formula retains the subscript 2 on the standard basis in three places. These should use n.
\item The diagram labels the downward domain change from B\_1 to B\_2 by Q, but the following equation defines Q as the change in the opposite direction, from B\_2 to B\_1.
\end{enumerate}
\section*{Projection: source notes}
\begin{enumerate}
\item The expanded first derivative omits a factor of two on v\_1s\_1+v\_2s\_2. The final stationary-point formula is correct, but the intermediate equality is not.
\item The second derivative is positive whenever the direction vector is nonzero, not only when neither component is zero. The one-zero-component case is not a trivial question.
\item The displayed calculation divides by v dot v although the answer also covers v=0. The zero case requires a separate nonzero line generator.
\item The closest-point comparison follows from the right angle and the Pythagorean theorem; the ordinary triangle inequality alone does not justify this hypotenuse comparison.
\item A comma is missing between the second and third vectors in this displayed ordered orthonormal basis.
\item The sentence repeats the word can.
\item The orthogonal-vector exercise permits zero vectors, but its projection formulas use each vector as a line generator and need nonzero denominators.
\item The last projection subscript uses v\_k instead of the bound orthogonal vector kappa\_k.
\item These answer formulas use v\_k or v\_2 where the corresponding bound orthogonal vector is kappa\_k or kappa\_2.
\item A displayed dot product takes a scalar fraction as its second argument before multiplying by kappa\_i; the dot product should instead enclose the entire scalar-vector product.
\item The orthogonal-basis calculation uses beta subscripts from the prior nonorthogonal basis while the explicit vectors are the current kappa basis.
\item The squared norm in the Bessel example is 30, not 50.
\item The first proof sentence invokes the null space of orthogonal projection before this lemma establishes that projection along M perp is well-defined. This is a proof-order dependency gap, not a false conclusion.
\item An equals sign is missing between B\_N and the displayed basis sequence.
\item The phrase all the the repeats the word the.
\item An equals sign is missing after M perp before its set description.
\item The follow-up projection answer uses the wrong plane matrix and incorrectly says (1,2,-1) belongs to S. Its correct orthogonal projection onto the stated S is (0,2,0).
\item As in the earlier line case, the closest-point proof needs the right-triangle Pythagorean comparison; the ordinary triangle inequality alone does not imply this conclusion.
\item An equals sign is missing between the named projection and its matrix-vector expression.
\item The final scalar coefficient is joined to w\_n by dotprod instead of scalar multiplication.
\item An equals sign is missing between the null-space symbol and its set description.
\item The sentence To finish we taking a basis has a verb-form error.
\end{enumerate}
\section*{Line of Best Fit: source notes}
\begin{enumerate}
\item The least-squares fit minimizes the sum of squared vertical residuals, not their total absolute length. In the printed coin data, m=79/140 gives squared cost 13/14 and absolute cost 9/7; m=17/30 gives squared cost 1 but the smaller absolute cost 1. The original wording and diagram are preserved.
\item The original sentence repeats “the” before “coefficients”. The unchanged original text is retained.
\item In supplied answer 7(a), the exact intercept 4259/1398 is correct but its printed decimal 3.239628 is not: it is approximately 3.046495. The seven-flight fit predicts 6317/2796, approximately 2.259299, failures at 31 degrees F (not 2.45); its four-failure threshold is -2666/71, approximately -37.549296 degrees F (not -29.94). These are mathematical checks of the historical classroom data, not safety guidance. Part (b) gives the correct all-flight coefficients; 2.79 is correctly rounded and the printed threshold 11 is a coarse rounding of 810/73, approximately 11.095890. Original formulas and answers remain unchanged.
\item The final supplied answer states orthogonality but cites the Triangle Inequality for a conclusion that follows from Pythagoras: squared distance to w equals squared distance to the projection plus squared distance from the projection to w. The Triangle Inequality alone does not supply that lower bound. The original “hypoteneuse” spelling and reasoning are preserved.
\item The source describes the asteroid belt as the remains of a planet that broke apart. NASA describes remnants from solar-system formation whose growth into planetary bodies was interrupted by Jupiter. Retain the exercise as a historical distance-fitting example, not as a current account of asteroid-belt formation; its original text is unchanged.
\item The source says this distance-fitting method helped discover Neptune. NASA credits the prediction to calculations of perturbations in Uranus’s orbit and the subsequent telescope search using Le Verrier’s position. Those dynamical calculations should not be conflated with this ordinal-distance least-squares exercise. The original historical assertion remains unchanged.
\item Question 8(d) requests Neptune’s predicted location, 8(e) requests Pluto’s, and 8(f) asks about accuracy for both. Supplied answer (d) gives the fitting coefficients only; answer (e) evaluates the Neptune index x=9 and compares its distance, and answer (f) evaluates Pluto at x=10 and compares its distance. The computations are present collectively but the last three answer labels do not match the corresponding tasks. Preserve the original labels and disclose this crosswalk separately.
\item The original two-column matrix in supplied answer 3 has a middle row with two successive vdots commands but no column separator. Other rows contain two cells. The original math is preserved rather than silently inserting an ampersand; this is a source-level layout defect, not a conversion-introduced value change.
\item The bridge-data description says “if a crossings” rather than the singular “if a crossing”. Its units and original wording are preserved.
\end{enumerate}
\section*{Geometry of Linear Maps: source notes}
\begin{enumerate}
\item Exercise 1 prints the row-reduction product as P and a Q that swaps columns 1 and 2, then claims H=PBQ. Exact rational replay disproves that product. The described column operations give B=P H (D S23); solving for H requires the inverses of those reduction products. Original matrices remain unchanged.
\item Both Exercise 2 rotation representations are juxtaposed with their matrices without equals signs. The mirror preserves those exact expressions; it does not silently insert a sign.
\item Exercise 3 describes the general counterclockwise angle theta but labels the representation t\_\{pi/4\}, while the matrix uses theta. Use a general theta label or explicitly specialize the matrix; this mirror preserves the original mismatch.
\item The next Exercise 3 representation writes t\_\{-pi/4\}, with literal letters p and i rather than the TeX command for pi. MathML conversion preserves this source error.
\item Exercise 6 evaluates the inverse derivative at x instead of at f(x). With differentiability and a nonzero derivative, the relation is (f\textasciicircum{}\{-1\})'(f(x))=1/f'(x). For the invertible real function f(x)=x+x\textasciicircum{}3, the inverse derivative at 2 is 1/4, not 1/f'(2)=1/13. The original formula remains unchanged.
\item The body states that any line maps to a line without mentioning collapse to a point. Exercise 5 and its answer explicitly allow a degenerate line and give the condition h(w)=0. Carry that qualification with modular extracts of the earlier narrative. This is a later source clarification, not a newly invented correction.
\end{enumerate}
\section*{Magic Squares: source notes}
\begin{enumerate}
\item The opening dimension formula says n²−n for magic squares when n≥3. The same active text gives dim M\_n=1+dim M\_n,0 and dim M\_n,0=n²−2n−1, which instead imply n²−2n. Exact rational rank replay gives dimension 3, not 6, at n=3, and agrees with n²−2n for n=3 through 12. The finite replay is not a new proof for all n; the original formula is preserved and this warning is separate.
\item The proof twice indexes inner columns by 2,…,n−2. The described first-block inner columns are 2,…,n−1: the source itself says column 2 for n=3 and columns 2 and 3 for n=4. The source coefficient arrays confirm those entries. Both occurrences are preserved unchanged.
\item Exercise 2 has three inconsistencies in its displayed reduction. The swap arrow has an extra leading minus although the next matrix is an ordinary row-2/row-6 swap. The later row-3 result requires +ρ₂+ρ₃, not the printed −ρ₂+ρ₃. The row-4 result of −ρ₂+ρ₄ must be (0,0,−1,1|0), not the displayed (0,1,−1,1|0). Exact symbolic replay verifies these discrepancies; a=b=c=d=s/2 does satisfy the original six equations. Original arrows and matrices are unchanged.
\item Exercise 4 defines θ(M) using Tr*(m), a lowercase variable not introduced there, while the supplied answer uses Tr*(M). The source argument should consistently refer to the matrix M. Both source versions remain unchanged.
\item Exercise 1’s median explanation says “shows that” immediately followed by an extra “Thus”. This is a duplicated prose connector, preserved in the original answer.
\end{enumerate}
\section*{Markov Chains: source notes}
\begin{enumerate}
\item The Sage matrix differs from the printed absorbing matrix at row 2/column 1 and row 5/column 6 (one-based): both code entries are .5 instead of 0. Transposition for the row-vector API does not repair those entries.
\item The listed p5 recurrence drops its absorbing-state carry term and does not supply the p4 interior recurrence used by the parity proof.
\item The four historical B100 commands use starts 1,1,3,1 dollars, while the immediately following four output columns are labelled starts 1,2,3,4 dollars.
\item The industrial matrix and NE-to-W prose use a row-oriented transition convention, but the answer multiplies that same matrix by column vectors. Its column sums are not 1; the first resulting purported probability vector has mass 1.323725. The final row’s .999 sum is a separate rounding residual, not the cause of the orientation error.
\end{enumerate}
\section*{Orthonormal Matrices: source notes}
\begin{enumerate}
\item The two quarter-turn directions in the original prose are reversed. For a=1 and b=0, the first displayed matrix is the identity, so its second basis vector is counterclockwise from its first. The second matrix reverses that orientation. The later orientation paragraph and original figures agree with the matrices; all original words remain unchanged.
\item Exercise 2(a) correctly uses y cos(pi/6) before simplification, but prints y cos(sqrt(3)/2) afterwards. Substitution should give the coefficient sqrt(3)/2. The original formula is preserved, and this separate note does not certify the rest of the answer.
\item Exercise 2(b) supplies a rotation, with determinant +1, rather than reflection in y=2x. Reflection has matrix [[-3/5,4/5],[4/5,3/5]]: it fixes (1,2), squares to the identity and has determinant -1. The original supplied formula is unchanged.
\item Exercise 2(c) again supplies a rotation rather than the requested reflection. The reflection in y=-2x has linear part [[-3/5,-4/5],[-4/5,3/5]], followed by translation (1,1). That translation does not change the determinant of the linear part. All original formulas remain unchanged.
\item The statement that any nontrivial translation following reflection gives a glide reflection needs a convention or qualification. Reflection in y=0 followed by translation (0,2) is just reflection in y=1 and fixes that line. A genuine nonzero glide requires a nonzero parallel component. This is a classification qualification, not a silent change to the source.
\item The similarity statement introduces points p and q but uses v in q=(kT)v+p\_0. This appears to be a p/v notation inconsistency; no variable is silently renamed in the original formula.
\end{enumerate}
\section*{Determinants: source notes}
\begin{enumerate}
\item The final elimination step in the supplied three-by-three determinant proof needs a minus sign: subtract ((ah-bg)/a) times row 2 from row 3. The printed plus sign does not give the displayed zero. The original operation and result remain unchanged.
\item In the a nonzero, ae-bd=0 case, the supplied answer reverses singular and nonsingular. Its zero-product conditions describe singularity; nonsingularity requires both remaining pivots to be nonzero. The later equality-to-zero conclusion has the same reversal. Original wording is retained.
\item The displayed two-by-two mnemonic expansion cancels identically to zero, yet the next line equates it with ad-bc. The identity matrix gives zero in that displayed expansion and determinant one. This concerns the printed expansion, not a claim that the ordinary two-by-two determinant formula is wrong.
\item The proposed function d(T)=det(TS)/det(S) is undefined for singular S. This proof route needs a nonzero determinant assumption and a separate singular-S argument. The product identity itself is true; the source exercise and answer are preserved.
\item Repeated interpolation abscissae make the Vandermonde system singular, but do not always make it inconsistent. Two copies of the point (0,1) are satisfied by every polynomial 1+ax; taking nonzero a gives infinitely many degree-one examples. Zero determinant means no unique solution, not necessarily no solution.
\item The supplied Vandermonde answer needs subtraction of x2 times row 2 from row 3. The printed addition does not produce the displayed triangular determinant. For x1=0,x2=1,x3=2, addition gives row (0,2,6), whereas subtraction gives (0,0,2). Original bytes remain unchanged.
\item The supplied answer incorrectly concludes that every reversed n-permutation has sign +1 for n greater than four. Its sign is (-1)\textasciicircum{}(n(n-1)/2), because every pair is an inversion. n=6 and n=7 both have sign -1. Applying the displayed parity pattern to n!, rather than n, is the faulty step.
\end{enumerate}
\section*{Geometry of Determinants: source notes}
\begin{enumerate}
\item The third volume answer gives -58, the signed determinant. The section defines volume as the absolute value, so the volume is 58. Original question/answer bytes are retained.
\item The supplied size factor for (x,y) maps to (3x-y,-2x+y) is -1, but its matrix determinant is +1. This transformation preserves oriented size. Original answer is retained.
\item The two orthogonality equations equate a scalar dot product to a three-component zero vector. Their right sides should be scalar zero. The subsequent linear system uses scalar zero and the final triangle area is consistent; no formula is silently changed.
\item The displayed ordered basis ((0,1),(-1,0)) is the standard basis rotated counterclockwise through pi/2, not the supplied pi/4. Its positive orientation conclusion remains correct. Original wording is retained.
\item The matrix-vector formula starts with t(s\_i) but uses the coordinates and expansion of s\_j. A single consistent column index is needed throughout. This is an original index inconsistency, not a failure of the determinant product theorem.
\item The question gives one half of the signed determinant as geometric area. Reversing vertex order makes that expression negative. Area requires its absolute value; the original question remains unchanged.
\item The positive-integer conclusion uses the nondegenerate meaning of triangle. If degenerate triples are admitted, collinear integer points give zero area. Carry that interpretation with the question; this is an assumption-scope note, not a confirmed error under the standard nondegenerate definition. Original question is retained.
\item The supplied slope/intercept formula does not follow from the preceding determinant equation. For points (1,0) and (3,2), it gives y=x+1, which passes through neither point; the correct line is y=x-1. Original expansion/formula are preserved.
\item The geometric explanation compares (x2,y3) with (x3,y3), accidentally excluding distinct vertical-line points. It should compare the actual points (x2,y2) and (x3,y3). Original condition is retained.
\item The prose says the illustrated shear uses k = -0.35. The original MetaPost construction uses z4 = -0.65 z1 and z5 = z4 + z2, so its pictured second vector is -0.65 v + w. Its adjacent comment says k = 0.35 and is inconsistent too. This records a disagreement between original witnesses; neither prose nor diagram is changed, and the intended parameter is not inferred.
\item The original diagram places v at (1,2) on its standard unit axes, while the adjacent coordinate column gives (3,2). Both are in the first quadrant, but they are different vectors. The exact original diagram and coordinate column are preserved. This note does not choose which witness the author intended and contains no supplied-answer reasoning.
\end{enumerate}
\section*{Laplace’s Formula: source notes}
\begin{enumerate}
\item The supplied Laplace proof names ex:ExpThreeFirstRow, but the frozen native label catalogue has no such target. Preserve the spelling and unresolved warning; do not guess a replacement example.
\item The inverse scalar must be 1/(t11*t22 - t12*t21), with the entire determinant grouped in the denominator and required to be nonzero. The printed 1/t11 t22 - t12 t21 lacks that grouping. For [[2,1],[1,1]] the required scalar is 1, not -1/2. The source formula is retained.
\item The definition says above the diagonal, but i > j denotes below it. Upper triangular matrices have zero entries below the diagonal. The displayed inequality is consistent with that meaning; the word above is not. This note concerns the question, not its solution.
\item The answer tries to prove M\_ab = 0 for a > b, although its own matrix has M21 = -12. To make adj(M) upper triangular one needs M\_ab = 0 for a < b. The following minor-index cases and the asserted zero entry above the diagonal are also reversed; changing one inequality alone does not repair this proof. For a < b, the first b-1 columns of the minor are supported in only b-2 surviving rows, so they are dependent and its determinant is zero. The theorem conclusion is correct, but the original argument is not preserved as a verified proof.
\item The definition uses T\_ij for a scalar cofactor, but this answer calls it a minor matrix and later takes det(T\_ij). Distinguish the deleted-row/column matrix M\_ij from the cofactor C\_ij = (-1)\textasciicircum{}(i+j) det(M\_ij). The permutation coefficient is C\_ij; for i = j = n its sign is +1. The original notation is retained, with this type/scope warning.
\item The column expansion of S = transpose(T) needs S\_ji = t\_ij, not t\_ji. With T = [[1,2],[3,4]], the printed coefficients give -5 while det(S) is -2. The final line also interchanges minor/cofactor notation and indices. The induction should compare the (j,i) deleted minor of S with the transpose of the (i,j) deleted minor of T, then equate the corresponding cofactors using their equal signs.
\item Read only n-2 nonzero entries as at most n-2: some diagonal entries may themselves be zero. With that interpretation the zero-row/column argument remains valid. This is a clarification of the bound, not a demonstrated error if only was intended to mean at most.
\item The preceding discussion explicitly proves T adj(T) = det(T) I by row expansion. The claimed identity adj(T) T = det(T) I needs the analogous column-expansion argument. A reusable proof must retain that second argument or identify it as the omitted analogue; the one-line pointer alone is not a complete standalone proof. No original theorem or proof is replaced.
\end{enumerate}
\section*{Cramer’s Rule: source notes}
\begin{enumerate}
\item The supplied answer claims that all replacement determinants vanishing characterizes infinitely many solutions. Without an additional rank hypothesis this is false: the zero 2-by-2 coefficient matrix has all replacement determinants zero both for b=(0,0), with infinitely many solutions, and for b=(1,0), with none. Cramer determinants alone do not distinguish these two singular systems. Preserve the original answer and attach this separate counterexample.
\item The phrase two nonsingular cases is inconsistent with the displayed coefficient matrix [[1,2],[1,2]], whose determinant is zero. Both the infinitely-many and no-solution examples are singular. The rest of the stated c=6 versus c!=6 distinction is consistent with the equations.
\end{enumerate}
\section*{Speed of Calculating Determinants: source notes}
\begin{enumerate}
\item The original plot uses 1.054 seconds for 50 rows, whereas the original transcript prints 1.05453586578 seconds, which rounds to 1.055 at three decimal places. Preserve both source values; this separate disclosure must accompany any presentation that claims the plot is the rounded transcript. No replacement timing experiment or altered original graph is warranted.
\end{enumerate}
\section*{Chiò’s Method: source notes}
\begin{enumerate}
\item In the initial elimination display (**), the final row should be (0,9,1): (2,10,2) minus (2,1,1). The later condensation matrix with final row (9,1) and determinant result 25 is correct. Preserve both printed displays and attach this separate correction; do not replace the later correct values.
\item The displayed (n-1)-by-(n-1) matrix C uses indices 1 <= i,j <= n-1, with entries a11*a(i+1,j+1)-a(i+1,1)*a(1,j+1). The printed range excludes the first row/column and runs past A.
\item The answer suggests any nonzero pivot but its displayed minor still uses row and column 1. Bring a chosen nonzero entry to (1,1) by specified row and column swaps, track the sign of each swap, then apply the ordinary formula to the permuted matrix. Return to the original determinant using that sign. If every entry is zero, the determinant is zero. Do not silently use the printed first-pivot formula for an arbitrary pivot.
\item The displayed remaining columns of B omit its first column (a11,0,...,0). For every i,j >= 2 the correct remaining entry is a11*aij-ai1*a1j. In particular, the last two printed second-row entries use ann and a2j where a11 and a21 are required. Retain the source display and read it with these exact qualifications.
\item The final comparison is between two expressions for det(B), not the printed set(B). The cancellation gives det(A)=det(C)/a11\textasciicircum{}(n-2), not det(C)/ann\textasciicircum{}(n-2), consistently with the topic's earlier statement. It requires the nonzero first pivot over the real-number field used here.
\end{enumerate}
\section*{Projective Geometry: source notes}
\begin{enumerate}
\item The frame lemma requires that no three of T1,U1,V1,O be collinear. The later normalization (t2,1,1) excludes T2=T1: the expression a(1,1,1)+b(1,0,0) can be divided by a only when a is nonzero. The U and V normalizations have the analogous restriction. Coincident side pairs also do not determine a unique intersection. These are scope requirements of the printed coordinate proof; no missing degenerate-case proof is supplied by this adapter.
\item The answer says any point on T0V0 has a representative (u0,1,1), but the endpoint T0=(1,0,0) does not: no nonzero scalar changes its two zero coordinates to ones. The frame and subsequent divisions therefore use a nondegenerate configuration. Retain the original wording and attach this explicit endpoint/scope qualification, rather than silently extending the answer to undefined joins or intersections.
\item The cross-product formula for the line through two points requires distinct projective points, hence nonproportional representatives; otherwise it returns the forbidden zero triple. Dually, two lines must be distinct. Preserve these elementary nondegeneracy conditions when reusing the exercise.
\end{enumerate}
\section*{Computer Graphics: source notes}
\begin{enumerate}
\item Division by c and the representative with third coordinate 1 apply only when c is nonzero. A nonzero vector with c=0 represents a projective point at infinity and has no representative on z=1. The origin of the affine screen is (0,0,1), not the excluded zero homogeneous vector.
\item In the last supplied answer, the z-axis matrix has a zero third row and the y-axis matrix has a zero second row. They erase the coordinate on the purported fixed axis, so neither is a rotation; even at theta=0 they are not the identity. The fixed-axis diagonal entries must be 1: entry (3,3) in the first matrix and (2,2) in the second. All other printed coefficients, including the y-rotation sign convention, are retained. This note does not replace the printed source.
\item An invertible projective transformation is represented by a nonsingular 3-by-3 matrix up to nonzero scale. Writing its bottom-right entry as 1 assumes that entry is nonzero; an invertible matrix with bottom-right entry 0 cannot be scaled to the displayed form. General affine maps may be singular; an affine automorphism instead requires an invertible 2-by-2 linear block. The printed matrices are unchanged.
\item Five displayed first components print cos without a leading backslash, unlike the matching cosine entries in their matrices. They refer to the cosine function, but the original TeX typesets three variable letters. Preserve the native mathematical bytes and attach this explicit operator-typography note.
\item The prose calls figure ch4.62 a rotation by half a radian. Its original MetaPost computes theta=0.5*(180/3.414159) degrees, using 3.414159 rather than pi. The rendered source figure is therefore approximately 0.460083 radians, not 0.5. Preserve this original figure and disclose the discrepancy; do not silently redraw it.
\end{enumerate}
\section*{Complex Vector Spaces: source notes}
\begin{enumerate}
\item Read this factorization statement for a nonzero polynomial with its leading scalar retained. Uniqueness requires normalized (monic) factors and ignores factor order; arbitrary nonzero rescalings of factors are otherwise possible. The zero polynomial does not have such a unique factorization. This is a separate qualification, not a silent change to the source.
\item For the factorization formulation, retain the nonzero leading scalar and factor order convention: a nonconstant polynomial factors into a leading scalar times monic linear factors, unique up to ordering. A nonzero constant has an empty linear-factor product; the zero polynomial is excluded from the uniqueness assertion. The source review cites an external full proof; this edition does not manufacture one.
\item The earlier native definition says degree n or less. That bound must remain in force here when changing the coefficients to complex numbers: the set of polynomials of exactly degree n is not the stated vector space. The original sentence remains untouched; the complete earlier definition is attached as evidence.
\item The preceding paragraph uses an informal definition based on factorization into two lower-degree polynomials. Preserve that local convention when quoting this sentence; do not export the assertion about constants as the usual ring-theoretic definition of an irreducible element. This is a convention boundary, not a silent replacement of the author’s definition.
\end{enumerate}
